O texto se apresenta como proposição experimental contra-objetivista. Desde o ponto de vista que propõe enxergar os diferentes corpos como sujeitos e/ou orientes (ver entradas abaixo) no lugar de uma perspectiva objetivadora perante o mundo e seus compartilhantes, escrevo na contramão da objetividade moderna. Assim sendo, o texto convida a uma comunicação que se dá mediante tropeços, incompletudes e proposições que se articulam, por vezes, contrariamente ao que esperaria de uma concordância objetivista.